% Generated draft: requires mathematical and source review.
\documentclass[11pt]{article}
\usepackage[margin=25mm]{geometry}
\usepackage{fontspec}
\setmainfont{Latin Modern Roman}
\usepackage{amsmath,amssymb,mathtools,unicode-math}
\setmathfont{Latin Modern Math}
\usepackage{enumitem,needspace,xurl,hyperref}
\setlength{\parindent}{0pt}
\setlength{\parskip}{5pt}
\setlength{\emergencystretch}{3em}
\newcommand{\N}{\mathbb{N}}
\newcommand{\Z}{\mathbb{Z}}
\newcommand{\Q}{\mathbb{Q}}
\newcommand{\R}{\mathbb{R}}
\newcommand{\C}{\mathbb{C}}
\newcommand{\F}{\mathbb{F}}
\newcommand{\A}{\mathbb{A}}
\newcommand{\PP}{\mathbb{P}}
\newcommand{\E}{\mathbb{E}}
\newcommand{\eps}{\varepsilon}
\newcommand{\dd}{\,\mathrm{d}}
\begin{document}
\section*{UnsolvedMath ID: 701038}
\textbf{Source problem code:} catalogue entry
\subsection{Definitions and mathematical statement}
A profinite group is a topological group isomorphic to an inverse limit of finite discrete groups. For a profinite group $G$, let $\mu_G$ denote its normalized Haar measure, so that $\mu_G(G)=1$. For a positive integer $n$, define the closed, hence measurable, subset
\[
X_n(G)=\{x\in G:x^n=1_G\}.
\]
\textbf{Conjecture (L. Levai and L. Pyber, 1999).} For every positive integer $n$ and every profinite group $G$, if $\mu_G(X_n(G))>0$, then there exist an open subgroup $H\leq G$ and an element $t\in G$ such that
\[
tH\subseteq X_n(G),\qquad\text{equivalently}\qquad (th)^n=1_G\quad\text{for every }h\in H.
\]
Thus every element of the left coset $tH=\{th:h\in H\}$ has finite order dividing $n$. No normality assumption on $H$ is imposed. The source leaves the range of $n$ implicit; the formulation above uses the conventional positive-integer interpretation.

\textbf{Supplied qualifications.} The source states that the conjecture holds for $n=2$. Its comment labelled 2021 states that it is also proved for $n=3$, citing a 2022 publication by A. Abdollahi and M. S. Malekan. Its comment labelled 2025 reports significant progress in the cited preprint by S. Kionke, N. Otmen, T. Toti, M. Vannacci, and T. Weigel, without specifying the precise results. These are supplied historical assessments, not independently verified claims. The registry labels the problem open.

The source also asks whether similar results hold when the solution set of some equation in a profinite group has positive measure. It does not specify the class of equations, the number of variables, whether coefficients are allowed, or the intended analogue of the coset conclusion; consequently this broader question does not determine a precise additional conjecture.
\subsection{Short English statement}
The conjecture says that if a positive Haar proportion of the elements of a profinite group satisfy $x^n=1$, then some coset of an open subgroup consists entirely of such elements. Here $n$ is interpreted as a positive integer. The supplied source reports the cases $n=2$ and $n=3$ as established and reports further progress in 2025, while the registry labels the general problem open; these assessments have not been independently verified.
\subsection{Sources}
\begin{itemize}
\item Local UnsolvedMath registry, permanent ID $701038$, exact source code \texttt{catalogue entry}, title \emph{Kourovka Notebook Problem 14.53}. The supplied record is labelled open and was updated on \texttt{2026-10-16T14:00:00Z}; that timestamp is registry metadata, not a source-verification date.
\item \emph{The Kourovka Notebook (Unsolved Problems in Group Theory)}, Problem 14.53, proposed by L. Levai and L. Pyber in the 14th issue (1999). The supplied background identifies the 21st issue edition (2026), arXiv:1401.0300v48. The supplied collection description names E. I. Khukhro and V. D. Mazurov as editors.
\item A. Abdollahi and M. S. Malekan, \emph{Adv. Group Theory Appl.}, \textbf{13} (2022), 71--81. Cited by the source's comment labelled 2021 for the case $n=3$; no article title or link is supplied.
\item S. Kionke, N. Otmen, T. Toti, M. Vannacci, and T. Weigel, preprint (2025), \url{https://arxiv.org/pdf/2507.19086}. The source's comment labelled 2025 attributes significant progress to this preprint without specifying its conclusions.
\end{itemize}
\end{document}
